%=============================================================================
% chktex-file 1 chktex-file 8 chktex-file 12 chktex-file 13 chktex-file 24 chktex-file 26
% CONCLUSION GÉNÉRALE ET PERSPECTIVES
%=============================================================================

\chapter*{Conclusion Générale et Perspectives}
\addcontentsline{toc}{chapter}{Conclusion Générale et Perspectives}
\markboth{Conclusion Générale et Perspectives}{Conclusion Générale et Perspectives}

\section*{Bilan du Projet et Réalisations}

Ce Projet de Fin d'Études, réalisé au sein de \textbf{\hostCompany}, avait pour ambition de concevoir et de développer \textbf{\projectShortTitle}, une plateforme logicielle SaaS B2B dédiée à la digitalisation intégrale du cycle de vie du recouvrement dans le secteur du crédit-bail financier.

Face aux limites critiques des méthodes traditionnelles (gestion sur feuilles Excel, risques de péremption légale, lenteur du traitement manuel des rapports d'expertise), notre travail a abouti à une solution logicielle industrielle articulée autour d'innovations majeures :
\begin{itemize}
    \item \textbf{Une architecture microservices polyglotte et sécurisée :} L'association de Next.js 16 (BFF avec cookies HttpOnly), Spring Boot 4 (Java 21) et FastAPI (Python) offre un équilibre optimal entre performance, maintenabilité et sécurité d'entreprise.
    \item \textbf{Une isolation multi-tenant par schémas SQL étanches :} L'étanchéité absolue des données entre institutions financières concurrentes garantit la conformité stricte avec les exigences réglementaires (RGPD et Loi n°63-2004).
    \item \textbf{Un workflow réglementaire en 5 phases :} Le guidage strict des étapes juridiques couplé à un moteur d'alertes préventives élimine les risques d'invalidation de procédure.
    \item \textbf{Un pipeline d'intelligence artificielle asynchrone :} L'extraction automatique des données des rapports d'expertise par NLP/LLM réduit le temps de traitement de 45 minutes à moins de 30 secondes par dossier, tout en fournissant une confrontation financière immédiate entre valeur d'expertise ($\ve$) et valeur résiduelle ($\vr$).
\end{itemize}

\section*{Apports Personnels et Professionnels}

Sur le plan professionnel et académique, ce projet a constitué une opportunité exceptionnelle d'aborder des problématiques d'ingénierie logicielle de pointe :
\begin{itemize}
    \item Maîtrise de la conception orientée domaine (\textit{Domain-Driven Design}) dans le secteur Fintech du crédit-bail.
    \item Implémentation pratique de patrons d'architecture SaaS multi-tenant complexes.
    \item Intégration de modèles de Traitement Automatique du Langage Naturel (NLP) et de LLM dans une chaîne d'inférence asynchrone en production.
    \item Application rigoureuse de la méthodologie Agile Scrum en autonomie complète de développement.
\end{itemize}

\section*{Perspectives d'Évolution et Travaux Futurs}

La plateforme LeasRecover offre un socle technologique évolutif ouvrant la voie à plusieurs axes de développement prometteurs pour les versions ultérieures (V2) :

\begin{enumerate}
    \item \textbf{Génération Automatisée d'Actes Légaux et Signature Électronique :} Intégrer un moteur d'édition dynamique de documents juridiques (mises en demeure, requêtes en saisie-revendication) avec signature électronique qualifiée conforme eIDAS.
    \item \textbf{Scoring Prédictif du Recouvrement :} Exploiter des algorithmes de \textit{Machine Learning} supervisé pour prédire la probabilité de recouvrement amiable dès les premiers jours d'impayé, afin d'optimiser l'allocation des ressources juridiques.
    \item \textbf{Connecteurs d'Intégration Bancaire (Open Banking / ERP) :} Développer des adaptateurs standardisés pour interconnecter nativement LeasRecover aux principaux progiciels de gestion de contrats (LMS / Core Banking) du marché.
\end{enumerate}
